\begingroup
\selectlanguage{english}
\noindent\textbf{Relative homotopy types, reflection, and integration.}

These pages occupy a hinge in the construction of $\HOT$. Grothendieck has accumulated several ways of representing a homotopy type over a small category $S$: a presheaf of categories, a split or unsplit fibration, a smooth category over $S$, or simply an arbitrary functor $X\to S$. He now wants to know exactly which localizations make these presentations interchangeable. Once that question is settled, the base may vary and the categories so obtained can begin to carry the operations of a derivator. The movement from a comparison at fixed $S$ to inverse and direct images explains both the long abstract digression and the care devoted to the two triangular identities of an adjunction.

The indispensable distinction is between forgetting weak equivalences in the total category and testing them relative to the base. For a map $u:X\to Y$ over $S$, write
\[
\begin{aligned}
u\in W_S&\quad\Longleftrightarrow\quad u\in W,\\
u\in W_S^{\mathrm d}&\quad\Longleftrightarrow\quad
 (s\backslash u)\in W\quad\hbox{for every }s\in S,\\
u\in W_S^{\mathrm g}&\quad\Longleftrightarrow\quad
 (u/s)\in W\quad\hbox{for every }s\in S.
\end{aligned}
\]
Here $s\backslash X=(s\downarrow X)$ and $X/s=(X\downarrow s)$ are comma categories, not strict fibres. The two fine classes retain respectively contravariant and covariant variation along arrows of $S$. On suitably smooth or proper representatives, the corresponding comma categories may be replaced by strict fibres. On homotopically locally constant representatives, the coarse and both fine tests coincide under the additional fibration hypothesis used in this part of the manuscript. The names of the fifteen categories in the large diagram serve this comparison; they do not describe fifteen different homotopy theories.

Two replacements organize the whole argument. The comma replacement $\Phi_SX$ is a split fibration with fibre $s\backslash X$; its natural map from $X$ is a right-fine equivalence. It therefore changes the representative while retaining the relative homotopy type. The path replacement $C_SX$ has locally constant fibres and is connected to $X$ by a coarse equivalence. It retains the total homotopy type but generally changes the relative one. The distinction is expressed by a reflective localization
\[
\rho_S:\HOT_W(S)\rightleftarrows\HOT_W^{\mathrm{lc}}(S):i_S,
\qquad \rho_S\dashv i_S,
\qquad \rho_Si_S\simeq\mathrm{id}.
\]
The unit $X\to C_SX$ need not be invertible in $\HOT_W(S)$. This is why the second equivalence in Corollary~2 on manuscript page~120, which offers $C_S$ as an alternative to $\Phi_S$ for computing the inverse image of an arbitrary relative type, must be corrected. Its $C_S$ expression computes the inverse image after locally constant reflection. The commentary on the first part of Chapter~VI already explains the basic distinction; the present commentary applies it precisely where the manuscript temporarily loses it.

The opening overlap with VI-a completes the adjunction by comparing $i_{C(X)}$ with $C(i_X)$. Concatenation of paths gives a common retraction, and hence equality after localization. There is also a formal reason: for a coaugmented endofunctor, invertibility of both iterated units forces their compatibility and identifies the essential image as a reflective subcategory. This familiar recognition principle is stated in higher categorical form by \citet[Proposition~5.2.7.4]{Lurie2009HTT}. Thus the path construction supplies particularly useful representatives, but its composition law is stronger data than the adjunction itself needs.

The abstract discussion beginning on manuscript page~85 isolates precisely this point. In a category $\mathcal M$ with weak equivalences $W$ and a class of fibrations, Grothendieck compares the localization of $\mathcal M/S$ with that of its full subcategory of fibrant objects. A functorial replacement gives the equivalence directly. A replacement obtained only by lifting properties requires a separate argument: two choices of lift are compared through a factorization of a relative diagonal, and become equal after localization. The manuscript proves this comparison, rather than merely assuming that choices disappear. Theorem--Scholie~2 then introduces a finer class $W_S^0\subseteq W_S$ agreeing with $W_S$ between fibrant objects and constructs the associated reflection. Our exposition writes out the naturality checks and the independence of choices, and explains why the functorial-replacement case already yields the adjunction without the additional path-composition hypothesis $A'$. This is an application of established deformation theory, not a new localization theorem; compare \citet[Theorem~2.2.11]{Riehl2014}.

There are several distinct uses of the word fibration here. A Grothendieck fibration describes transport of objects in a category. A homotopical fibration controls comma categories and strict fibres. A fibration in a Quillen model structure is part of a lifting and factorization system. They interact, but the manuscript's abstract argument does not identify them. In particular, invoking Thomason's model structure gives a method for ordinary nerve weak equivalences on categories; it does not by itself establish every assertion for every proposed fundamental localizer. Likewise, weak saturation, meaning the formal closure axioms initially imposed on $W$, must be distinguished from strong saturation, which detects membership in $W$ by invertibility in $\mathrm{Cat}[W^{-1}]$. The latter is supplied for fundamental localizers by \citet[Proposition~4.2.4 in the cited author version]{Cisinski2006Presheaves}. It closes the explicit saturation reservation in the converse of Proposition~2 on manuscript pages~124--126.

Once the base varies, raw pullback of an arbitrary $X\to S$ is insufficient. To calculate $f^*$ one first takes a representative, such as $\Phi_SX$, on which strict pullback respects the relevant weak equivalences. Smoothness of $X\to S$ supplies one sufficient condition; properness of the change of base supplies another. The ``sure conjecture'' on manuscript page~122 is not left open in the work: Grothendieck proves it on manuscript pages~168--169 of the same chapter. The commentary separates this later proof from what has actually been established in these pages. For ordinary weak homotopy equivalences, testing the comma replacement of a point also proves the converse: preservation of right-fine equivalences by strict pullback characterizes properness. A three-object span gives the nonproper homotopy fibration requested on manuscript page~126; it preserves coarse equivalences but fails to preserve a particular fine equivalence. These computations make the difference between the two base-change conditions explicit.

A natural transformation $f\Rightarrow g$ also requires attention: it induces $g^*\Rightarrow f^*$ in the presheaf convention, but a literal functor between arbitrary strict fibres need not exist. Split fibrations make that action explicit; localization then removes the dependence on the chosen splitting.

The questions about the locally constant theory lead to a useful distinction in the modern answer. Restriction between categories of locally constant diagrams has both adjoints; this holds for arbitrary derivator coefficients by \citet[Corollary~3.11]{Cisinski2009LocallyConstant}. Nevertheless, the varying-base assignment $S\mapsto\HOT^{\mathrm{lc}}(S)$ does not form a derivator for ordinary homotopy types. The pointwise Kan extension formula already fails for a point included in the walking arrow. Fixing a base $S$ and varying an additional diagram category is a different operation: that construction does give the derivator of locally constant coefficients, as established in the same paper's Theorem~3.13. The distinction is visible without any higher machinery and prevents an attractive but incorrect explanation of Grothendieck's hesitation on manuscript page~128.

Another reservation concerns infinite disjoint unions of indexing categories. The corresponding product statement for arbitrary localizations is false. The bounded-length obstruction that Grothendieck notices can be realized by a countable family of finite fence categories: each factor becomes the terminal groupoid when all arrows are inverted, while the product has objects that no finite zigzag can connect. A refinement using finite categories inside a product with the walking retraction also answers footnote~66: strong saturation alone does not ensure that evaluation detects isomorphisms in localized diagram categories. A uniform three-arrow morphism is invertible at every coordinate, but its coordinate inverses require unbounded word lengths. The diagram shape in this counterexample is countably discrete; no finite-shape conclusion is asserted. The space derivator nevertheless satisfies both pointwise detection and the unrestricted disjoint-union axiom. Model structures and controlled fraction presentations supply the additional uniformity that the abstract questions lack.

The final pages construct integration. Postcomposition with $f:S'\to S$ sends $X'\to S'$ to the same total category over $S$, and preserves the fine weak equivalences. This yields $f_!$ before any choice of fibrant replacement. The unit uses the comma replacement over the new base; the counit for arbitrary $X$ is the localized composite
\[
\Phi_SX\times_S S'\longrightarrow\Phi_SX
  \xrightarrow{\,i_X^{-1}\,}X.
\]
The order matters. Pulling back $i_X$ and then trying to invert it can fail, because strict pullback need not preserve a fine equivalence between arbitrary representatives. Grothendieck notices and repairs that attempt on manuscript pages~142--143. At manuscript page~145 the argument has reached the last square in the second triangle; the commentary completes its verification from naturality. In the case of spaces this integration is homotopy left Kan extension, and for the structural map it recovers homotopy colimit. The present derivator formalism provides its axiomatic setting; see \citet{Groth2013}.

The modern picture may be expressed as coherent diagrams of spaces on $S^{\mathrm{op}}$, with locally constant objects corresponding to diagrams on the homotopy type of $S$. It is the entire homotopy type, not just the ordinary fundamental groupoid together with a functor into the homotopy category, that retains all the coherence. In particular, factoring the underlying diagram through $\Pi S$ detects local constancy but does not reconstruct its coherent object. The comma replacements and their categories of elements retain the additional data needed to form the homotopy colimit and to calculate its maps.
\endgroup
